\documentclass[11pt]{article}
\usepackage[margin=1in]{geometry}
\usepackage{amsmath,amssymb,amsthm,bm,mathtools}
\usepackage{algorithm}
\usepackage{algpseudocode}
\usepackage{enumitem}
\usepackage{hyperref}

% \documentclass{article}
\usepackage{graphicx} % Required for inserting images
\usepackage{hyperref}
% \usepackage[a4paper, total={5in, 10in}]{geometry}
% \usepackage{xcolor}
% \usepackage{color,soul}

% Math
\usepackage{amsmath}
\usepackage{amssymb}

% Theorems and definitions
\usepackage{amsthm}
\usepackage[english]{babel}

\theoremstyle{definition}
\newtheorem{definition}{Definition}[section]

\theoremstyle{remark}
\newtheorem*{remark}{Remark}


% Subplots
\usepackage{graphicx}
\usepackage{subcaption}

% colored area
\usepackage{tcolorbox}
\usepackage{xcolor} % Required for defining custom colors
\definecolor{lightgray}{gray}{0.9} % Define a light gray color

% Highlights
\usepackage{color,soul}

% Derivative
\usepackage{derivative}

\hypersetup{colorlinks=true,linkcolor=blue,citecolor=blue,urlcolor=blue}


\title{On the tradeoffs between generalized self-concordant convex functions and convex constraints}
% \author{(Draft)}
\date{\today}

% theorem envs
\theoremstyle{plain}
\newtheorem{theorem}{Theorem}
\newtheorem{proposition}{Proposition}
\theoremstyle{definition}
% \newtheorem{definition}{Definition}
% \newtheorem{assumption}{Assumption}
% \newtheorem{remark}{Remark}

\newcommand{\R}{\mathbb{R}}
\newcommand{\E}{\mathbb{E}}
\newcommand{\prox}{\mathrm{prox}}
\newcommand{\argmin}{\mathop{\mathrm{arg\,min}}}
\newcommand{\norm}[1]{\left\lVert #1 \right\rVert}
\newcommand{\ip}[2]{\left\langle #1,#2 \right\rangle}
\newcommand{\1}{\mathbf{1}}



% Set commands
\newcommand{\Y}{\mathcal{Y}}
\newcommand{\X}{\mathcal{X}}
\newcommand{\V}{\mathcal{V}}
\newcommand{\U}{\mathcal{U}}
\newcommand{\G}{\mathcal{G}}
\newcommand{\D}{\mathcal{D}}
\newcommand{\N}{\mathcal{N}}
\newcommand{\F}{\mathcal{F}}
\newcommand{\A}{\mathcal{A}}
% \newcommand{\T}{\mathcal{T}}
\renewcommand{\L}{\mathcal{L}}
\renewcommand{\P}{\mathcal{P}}

\renewcommand{\t}{^{\top}}
\newcommand{\qt}[1]{``#1''}
\newcommand{\citeReq}{{\color{red} (cite)}}

\newcommand{\Expected}[1]{\mathbb{E}[#1]}

% Math sets commands
\renewcommand{\R}{\mathbb{R}}


\begin{document}
\maketitle

% \begin{abstract}
% We present a rigorous optimization framework for fairness-aware regression
% where fairness is quantified through the \emph{pairwise mean residual gap}:
% the maximum absolute difference between average residuals across demographic groups.
% This fairness functional induces a nonsmooth, piecewise-linear penalty.
% We derive the full optimization formulation, algorithms, and trade-off analysis:
% proximal and primal–dual methods for solving the problem,
% envelope identities connecting fairness and accuracy,
% and sensitivity bounds describing how solutions evolve with the regularization weight~$\lambda$.
% \end{abstract}


\section*{Motivation}

There exists plenty of literature on how to solve convex optimization problems {\color{red} (citation)}. The same applies to the behavior of these problems when we want to modify in some way the feasible region of the problem, and want to find conditions on how our construction can still produce an optimal solution for the new problem or how this change would affect the optimal cost; i.e., sensitivity analysis {\color{red} (citation)}. In this work, we aim to mix both points of view to tackle the latter research problem, i.e, using the knowledge and techniques that are commonly seen in convergence of iterative algorithms (e.g. interior-point methods, Newton methods) to approximate the \textit{tradeoffs} between objective functions and changes in the (convex) feasible region. We focus on cases where relying on theoretical results for sensitivity analysis such as the \textit{envelope theorem} {(\color{red} cite)}, which depends on the optimal dual solution, is not possible since re-solving the problem for each desired change on the feasible region has a extremely high computational cost.  



Formally, we will consider settings where the objective function $f:\R^p \rightarrow \R$ is convex generalized self-concordant (e.g. mean squared error, entropy functions, log-barriers etc.) ---a concept introduced to analyze the convergence of Newton-like methods ({\color{red} cite})---, and where the feasible region is convex on every scenario. That is, we consider the convex optimization problem \begin{align}
    \begin{array}{crlc}
    \label{opt:P_tau}
      \P(\tau):=& \min_{x}  & f(x) \\
       % & s.t. &  x\in \X, \\
       & s.t. & h(x) \leq \tau, &
    \end{array}
\end{align}
where $h:\R^p \rightarrow \R$ is a convex function, and $\tau$ is a scalar that defines some level of constraint that we aim for. We further assume that we have access to a baseline optimal solution $x_0 \in \arg\min_{x\in \R^p} f(x)$. The quantity we are interest in studying is $\odv{f(x_\tau)}/{h}$, where $x_\tau$ is an optimal solution of $\P(\tau)$; which captures the exact \textit{tradeoff} between the functions $f$ and $h$ at the point $x_\tau$. Note that, for all $\tau \in \R$ such that $\P(\tau)$ is feasible, there exists an scalar $\lambda_\tau \geq 0$ such that $\P(\tau) = \min_{x\in \R^p} f(x) + \lambda_{\tau} h(x)$, and thus, using the \textit{envelope theorem}, we have that  
\begin{align*}
    \odv{f(x_\tau)}{h}=-\lambda_\tau.
\end{align*}
However, solving $\P(\tau)$ for each possible value of $\tau$ requires solving $\P(\tau)$ at least $\mathcal{O}\left({\log_2(\lambda_{{T}}-\lambda_{{0}})}/{\varepsilon}\right)$ times for $\lambda_\tau \in [\lambda_0, \lambda_T]$. Moreover, the structure of the each problem $\P(\tau)$ can be already highly expensive because of multiple reasons such as the nonlinear relationships, dimensionality of $\R^p$, number of constraints implicitly stated with $h$, etc. 

Begin able to efficiently compute the quantity $\odv{f(x_\tau)}/{h}$ ---or a close approximation for it--- is highly relevant for multiple applications such as constrained machine learning (regularization, fairness, etc.), resource allocation problems, portfolio optimization, $\ldots$ \citeReq. Thus, in this work we aim to, efficiently compute an approximation (exact derivation in some cases) for this quantity, which can be used to either make informed decisions on a neighborhood of $x_0$, or to even iteratively get an $\varepsilon$-approximation for it.

\section{Introduction}

Our approach is based on multiple assumptions that are satisfied by a very broad class of problems. We aim to approximate the optimal value for $f(x_\tau)$, given one optimal solution $x_\tau$ for problem (\ref{opt:P_tau}), using the threshold $\tau$ on $h(x_\tau)$. We assume that both $f$ and $h$ are convex functions on $\R^p$, with $f$ being self-concordant and possibly either strongly convex or smooth relative to some other function \citeReq. For $h$ we only require convexity, but tighter approximations can be achieved when assuming the same properties as $f$. We introduce these concepts in the following.

\subsection{Relative strong convexity and relative smoothness}

Here we introduce the generalization of the concept of $\mu$-strongly convex function and $L$-smooth function, introduced by \citeReq. 

\begin{definition}[Bregman divergence]
% Bregman divergence generated by a differentiable convex function \phi
\[
D_{\phi}(x,y)
:= \phi(x) - \phi(y) - \langle \nabla \phi(y),\, x - y \rangle ,
\]
where \(\phi:\mathbb{R}^n\to\mathbb{R}\) is differentiable and convex.
\end{definition}

\begin{definition}[Relative strong convexity]
$f$ is $\mu$-strongly convex relative to a reference function $\phi$ if 
\[
f(x) \;\ge\;
f(y) + \langle \nabla f(y),\, x - y \rangle
+ \mu\, D_{\phi}(x,y),
\qquad \forall x,y ,
\]
where \(D_{\phi}\) is the Bregman divergence of a strictly convex differentiable reference function \(\phi\).
\end{definition}

\begin{definition}[Relative smoothness]
 $f$ is $L$-smooth relative to a reference function $\phi$
\[
f(x) \;\le\;
f(y) + \langle \nabla f(y),\, x - y \rangle
+ L\, D_{\phi}(x,y),
\qquad \forall x,y ,
\]
where \(D_{\phi}\) is the Bregman divergence generated by a strictly convex differentiable reference function \(\phi\).
\end{definition}

\subsection{Generalized self-concordance}


The generalization of self-concordance proposed in \citeReq, which aimed to broaden the family of functions that could be classified as self-concordant, in order to derive efficient Newton-type methods for them. 
% \begin{definition}[Generalized self-concordance]
%     % f : (a,b) -> R is three times differentiable
% \[
% \text{A three-times differentiable convex function } \varphi:\R\to\mathbb{R}
% \text{ is called } (M_\varphi,v)\text{-generalized self-concordant if}
% \]
% \[
% \bigl|\varphi'''(t)\bigr|
% \;\le\;
% M_\varphi \,\bigl(\varphi''(t)\bigr)^{\frac{v}{2}},
% \qquad \forall t \in \textbf{dom}(\varphi) ,
% \]
% where \(M_\varphi \geq 0\) and \(v>0\).

% \end{definition}

\begin{definition}[Generalized self-concordance]
    A three times differentiable function $\varphi:\R \rightarrow \R$ is $(M_\varphi,v)\text{-generalized self-concordant if}$ \begin{equation}
        |\varphi'''(t)|\leq M_\varphi \varphi''(t)^{\frac{v}{2}}, \quad  \forall t \in \textbf{dom}(\varphi)
    \end{equation}
    where $M_\varphi\geq 0$ and $v>0$.
\end{definition}
As with the usual self-concordance definition, this concept can be extended to multivariate domain functions moving along the $x+td$ line $\psi(t)=f(x + td)$, but the equivalent definition is case-dependent, so we refer to \citeReq to get a proper understanding of the concept.

% \[
% \bigl|\nabla^3 f(x)[d,d,d]\bigr|
% \;\le\;
% \Bigl(\sum_{i=1}^k M_i\Bigr)\,\|d\|_{x}^{\,3-v}\,\|d\|^{\,v},
% \quad \forall x\in\mathbb{R}^n,\ \forall d\in\mathbb{R}^n.
% \]

\subsubsection{Some properties of self-concordance functions}

\begin{itemize}
    \item \textbf{Sum of self-concordant functions:} Let \(f_i:\mathbb{R}^n\to\mathbb{R}\), \(i=1,\dots,k\), be
\((M_i,v)\)-generalized self-concordant (with the same \(v\ge 2\)). Let $\beta\in \R^n$ such that $\beta_i>0, \forall i \in [n]$.
Then the sum
\[
f(x) := \sum_{i=1}^k \beta_if_i(x)
\]
is $(M_f,v)$-generalized self-concordant with the same $v\geq 2$ and $M_f:= \max\{\beta_i^{1-\frac{1}{v}}M_i \;|\; 1\leq i \leq m\}$.

\item \textbf{Bounds on self-concordant functions:} \begin{itemize}
    \item  Case $v=2$ {\color{red} (what we had)}. Let \(f:\mathbb{R}^n\to\mathbb{R}\) be \((M,2)\)-generalized self-concordant.
Fix \(x\in\mathbb{R}^n\) and direction \(d\in\mathbb{R}^n\), and define
\[
\phi(t) := f(x + t d), 
\qquad
\|d\|_x := \sqrt{d^\top \nabla^2 f(x)\, d},
\qquad
z := M\,|t|\,\|d\|_x.
\]
Then, for all \(t\) such that \(x + t d \in \mathrm{dom}(f)\) and \(z>0\),
% \[
% \phi(t) - \phi(0) - t\,\phi'(0)
% \;\ge\;
% \frac{e^{-z} + z - 1}{z^2}\,\|d\|_x^2,
% \]
% \[
% \frac{e^{-z} + z - 1}{z^2}\,\|d\|_x^2 \leq \phi(t) - \phi(0) - t\,\phi'(0)
% \;\le\;
% \frac{e^{\,z} - z - 1}{z^2}\,\|d\|_x^2.
% \]
% For \(z=0\), note that the limit as \(z\to 0\) gives
% \(\tfrac{1}{2}\|d\|_x^2\).

% - Case $v > 2$ {\color{red} (from the paper)}:  Let \(f:\mathbb{R}^n\to\mathbb{R}\) be \((M,v)\)-generalized self-concordant, with $v>2$.
% Fix \(x\in\mathbb{R}^n\) and direction \(d\in\mathbb{R}^n\), and define
% \[
% \phi(t) := f(x + t d), 
% \qquad
% \|d\|_x := \sqrt{d^\top \nabla^2 f(x)\, d},
% \qquad
% z := M\,|t|\,\|d\|_x.
% \]
% Then, for all \(t\) such that \(x + t d \in \mathrm{dom}(f)\) and \(z>0\),
% \[
% \phi(t) - \phi(0) - t\,\phi'(0)
% \;\ge\;
% \frac{e^{-z} + z - 1}{z^2}\,\|d\|_x^2,
% \]
\begin{equation}
\label{eq:GSC_exp_bounds}
    \frac{e^{-z} + z - 1}{z^2}\,\|d\|_x^2 \leq \phi(t) - \phi(0) - t\,\phi'(0)
\;\le\;
\frac{e^{\,z} - z - 1}{z^2}\,\|d\|_x^2.
\end{equation}
% \[

% \]
For \(z=0\), note that the limit as \(z\to 0\) gives
\(\tfrac{1}{2}\|d\|_x^2\).

\item Case $v>2$ {\color{red} (from paper): Much more complicated form}. 

% \item % Logarithmic lower/upper bounds for v=3 (standard self-concordant case)
% Let \(f:\mathbb{R}^n\to\mathbb{R}\) be a standard self-concordant function.
% Fix \(x\in\mathbb{R}^n\) and a direction \(d\in\mathbb{R}^n\), and define
% \[
% \phi(t) := f(x + t d),
% \qquad
% \|d\|_x := \sqrt{d^\top \nabla^2 f(x)\, d},
% \qquad
% \rho := |t|\,\|d\|_x.
% \]
% Then, for all \(t\) such that \(\rho < 1\), we have the two-sided bounds
% \[
% f(x + t d)
% \;\ge\;
% f(x) + t\,\langle \nabla f(x), d\rangle
% \;+\;
% \omega(\rho),
% \]
% \[
% f(x + t d)
% \;\le\;
% f(x) + t\,\langle \nabla f(x), d\rangle
% \;+\;
% \omega_{*}(\rho),
% \]
% where
% \[
% \omega(\rho) := \rho - \ln(1+\rho),
% \qquad
% \omega_{*}(\rho) := -\rho - \ln(1-\rho),
% \qquad 0 \le \rho < 1.
% \]


\end{itemize}

\end{itemize}


{\color{red} \textit{Note:} For now, we have only focused on the bounds of sums of self-concordant functions composed with affine functions that are $(M,2)$-generalized self-concordant. I think that should be enough for now, since it encompasses all the problems that we have explored and more: GLMs, SVM, portfolio optimization, some allocation problems, etc. This is also makes the bound on self-concordant functions straightforward and without dependence on $\norm{d}$ on the exponential or harder terms instead of the exponential ($v>2$ case).}

\section{Methodology}

The methodology we use for deriving bounds/approximations of the \textit{tradeoff} $\odv{f(x_\tau)}/{h}$ consists of two major steps, assuming $f$ is $(M_f, 2)$-generalized self-concordant and $h$ is either $(M_f, 2)$-generalized self-concordant or it is convex and its subgradient $\partial h$ is the combination of finite gradients of GSC functions:

\begin{enumerate}
    \item We solve the Taylor approximation problem of $\P(\tau)$ using the second-order approximation for $f$ and the first-order approximation of $h$ for the constraint, to get a closed-form efficient direction $d$ to move from $x_0$. That is, we solve  
    \begin{equation}
\label{opt:Taylor}
    \begin{array}{rlc}
      \displaystyle T(\tau):= \min_{d \in \R^p} &  f(x_0) + \frac{1}{2}\norm{d}^2_{x_0} \\
        s.t. & h(x_0) + {g_0} \t d \leq \tau, & %\\
        % & \beta \in  \beta
    \end{array}
\end{equation}
    where $g_0\in \partial h(x_0)$ is a subgradient vector of $h$ at $x_0$. Denoting $H_0 = \nabla^2f(x_0)$, this problem has a closed-form solution of the form $d^*= -\delta H_0^{-1}g_0 / \norm{g_0}_{H_0^{-1}}^2$, where $\delta = \tau - h(x_0) > 0$ (if not, there is no change). Further, if $\partial h(x_0)$ consists of the finite conic combination of the $m$ nonzero vectors $\{a_0^1, \ldots, a_0^m\}$, then we do not have to choose one subgradient, and can directly compute the solution as $$d^*= -\delta H_0^{-1}A_0(A_0\t H_0^{-1}A_0)^{-1}e,$$ where $A_0 = (a_0^1, \ldots, a_0^m) \in \R^{p \times m}$, and $e\R^m$ is a vector of ones.
    % Note that  if we use some constant $C_f$ that secured a 
    
    Since $f$ is $(M_f, 2)$-generalized self-concordant, this further leaves us with a baseline lower bound using the one explained in the previous section.

    \item Solving one extra 1-dimensional convex problem for each of the desired bounds. This is either upper/lower bound, or both, using the inequalities on (\ref{eq:GSC_exp_bounds}). In other words, we solve
        \begin{equation}
\label{opt:LB_approx}
    \begin{array}{rlc}
      \displaystyle \P_{LB}(\tau):= \min_{t \in \R_{\geq 0}} &  f(x_0) + \norm{d}^2_{x_0} \frac{(e^{tM_f } - tM_f -1)}{M_f^2 }\\
        s.t. & h(x_0) + {a_0^j} \t d +  \norm{d}^2_{H_0^j} \frac{(e^{tM_j } - tM_j -1)}{M_j^2 }\leq \tau, & \forall j \in [m],\\
        % & \beta \in  \beta
    \end{array}
\end{equation}
\begin{equation}
        \label{opt:UB_approx}
    \begin{array}{rlc}
      \displaystyle \P_{UB}(\tau):= \min_{t \in \R_{\geq 0}} &  f(x_0) + \norm{d}^2_{x_0} \frac{(e^{-tM_f } + tM_f -1)}{M_f^2 }\\
        s.t. & h(x_0) + {a_0^j} \t d +  \norm{d}^2_{H_0^j} \frac{(e^{-tM_j } + tM_j -1)}{M_j^2 }\leq \tau, & \forall j \in [m],%\\
        % & \beta \in  \beta
    \end{array}
\end{equation}
where $H_0^j$ are the hessians of each of the $m$ functions that define $h$. The \textit{tradeoff} then satisfies
\begin{align*}
    \odv{f(x_\tau^{LB})}{h}\leq \odv{f(x_\tau)}{h} \leq \odv{f(x_\tau^{UB})}{h}, \quad \forall \tau.
\end{align*}
\end{enumerate}


Note that if we instead of using the self-concordant assumptions, we make use of relative strong convexity and relative L-smoothness, we would get looser constraints, but we would just need to solve a quadratic equation instead of the convex problems.
% Assuming that $f:\R^n \rightarrow \R$ is a sum of self-concordant function (also self-concordant), then 

% % Multivariate generalized self-concordance
% \[
% f \text{ is } (M,v)\text{-generalized self-concordant if for all } x\in\mathbb{R}^n
% \text{ and all directions } d\in\mathbb{R}^n,
% \]
% \[
% \bigl|\nabla^{3} f(x)[d,d,d]\bigr|
% \;\le\;
% M \,\|d\|_{x}^{\,3-v}\, \|d\|^{\,v},
% \]
% where
% \[
% \|d\|_{x} := \sqrt{ d^\top \nabla^{2} f(x) d }
% \]
% is the local norm induced by the Hessian, and \(v\ge 2\), \(M>0\).




\end{document}
